% Appendix A, Small Molecules, hERG blockers held at separate sites -- c8-P7: secondary count no result depends on
For 56 of the 166 and 57 of the 170 shared molecules, Lhasa's call disagrees with the call from p$K_i$.

% Appendix A, Chapter 3 versions, Table tab:curation_ch3_ch5 -- c8-P7: scaffold-group counts (Chapter 5 values remain in Table tab:datasets)
Scaffold groups: LogD, Caco-2, hERG $K_i$ & 2,245; 1,038; 749 & 2,413; 1,086; 769 \\

% Appendix A, Chapter 3 versions, caption of tab:curation_ch3_ch5 -- c8-P7: caption sentence for the removed row
The last two rows give the numbers of molecules and of Bemis--Murcko scaffold groups.

% Appendix A, Antibodies, Nanobody thermostability -- c8-P7: measurement-method breakdown
The labels pool measurements from four main methods.

% Appendix A, Antibodies, Humanness -- c8-P7: species breakdown (removed phrase: 'among them humanized, chimeric and mouse antibodies,')
The other 366, among them humanized, chimeric and mouse antibodies, are labeled non-human (0).

% Appendix B, Data and code availability, opening paragraph -- c8-P3: QM9 pipeline is introduced again in the qsar_qm_models paragraph
Three separate pieces of software lie behind this thesis: NoiseInject, the label-noise package released with Chapter~\ref{chap:robustness}; a pipeline for that chapter's QM9 experiments; and the KIRBy repository,

% Appendix B, Data and code availability, NoiseInject paragraph -- c8-P2 option B: Chapter 2 already describes the package; kept the license and the pointer
% [TRIM2 2026-09-30] metric name: noise degradation slope -> normalized area under the retention curve (as in Chapter 2)
The NoiseInject benchmarking framework, implementing all six regression and six classification noise injection strategies together with the robustness metrics (normalized area under the retention curve and retention) and uncertainty-calibration metrics (ECE, coverage at $1\sigma$/$2\sigma$, mean interval width, and uncertainty--error / uncertainty--noise correlation) described in this work, is available as an open-source Python package under an MIT license \citep{noiseinject}. It also contains optional reference wrappers (a Gaussian process, split-conformal prediction, and Monte Carlo dropout) that produce per-sample uncertainty for use with these metrics. NoiseInject is model-agnostic, operates on label arrays (NumPy arrays or pandas Series), and is designed to integrate into existing cheminformatics or other data pipelines. Worked examples covering molecular-descriptor, fingerprint, and GNN-embedding workflows are provided in the repository.


% Appendix B, Data and code availability, hERG blocker paragraph -- c8-P4: counts are given in Appendix A (Table tab:curation_counts and the hERG blocker paragraph)
The exception is Lhasa Limited's curated hERG training set of 7,772 molecules and test set of 4,383 molecules \citep{Hanser2025},
